\section{Mở rộng sang Tiếng Anh và giọng nói}\label{sec:english-design}
Môn Tiếng Anh được thêm vào sau khi gia sư Toán đã được đo và duyệt, nên yêu cầu thiết kế quan trọng nhất là không làm đổi hành vi, công thức, ngưỡng, lời nhắc hay kết quả đánh giá của Toán và không để nội dung hai môn lẫn vào nhau. Hệ thống giữ điều này bằng cách đặt môn học thành thuộc tính bắt buộc ở mọi tầng thay vì cờ ở giao diện. Mọi tài liệu giáo viên, câu trong ngân hàng, gói bài và gói chương mang trường môn học, và thiếu trường này là lỗi kiểm tra dữ liệu. Bộ đọc tài liệu Tiếng Anh là mô-đun riêng được chọn theo môn của tài liệu mà không sửa nhánh xử lý của Toán. Mỗi môn có chỉ mục truy xuất dựng riêng từ đầu và bộ tách từ riêng, thay vì một chỉ mục chung rồi lọc lúc truy xuất. Môn của một buổi được ghi vào chính tài liệu buổi khi tạo, nên buổi cũ luôn được hiểu đúng môn. Các phần dùng chung, như phép chiếu mức hiểu, chính sách trên đồ thị, vòng dạy và cổng AI, chỉ được thêm một nhánh theo môn đọc từ gói nội dung của buổi, không tham số nào do máy khách chọn. Kiểm tra dữ liệu từ chối gói thiếu môn, kiểm tra dựng chỉ mục chứng minh chỉ mục hai môn không chứa đoạn của nhau, và bộ kiểm tra chuẩn của Toán được chạy lại trước mỗi thay đổi chạm phần dùng chung với yêu cầu kết quả giống hệt.

Phạm vi Tiếng Anh là Unit~11 ``Family time'' của lớp 5, gồm sáu bài về hỏi đáp chuyện đã xảy ra bằng thì quá khứ đơn (``Did you\ldots''), từ vựng về một chuyến đi và một đoạn đọc ngắn. Gia sư ôn bài theo cùng vòng sáu bước và cùng phép chiếu mức hiểu; cô giảng bằng tiếng Việt, còn câu đích học sinh nói hay chọn là tiếng Anh. Đáp án và câu hỏi chỉ lấy từ tài liệu chủ dự án cung cấp, và không mô hình nào được tạo đáp án cho câu kiểm tra. Bộ chấm nhận nhiều biến thể hợp lệ, phân biệt lỗi chính tả với lỗi ngữ pháp, coi một từ đúng là đúng một phần và nhận câu trả lời bằng tiếng Việt. Chính sách mức hiểu riêng có ba khác biệt: một lớp lỗi về ảnh hưởng của tiếng mẹ đẻ, câu nói theo cô được tính độc lập tối đa như bậc H6, và độ tin cậy nhận dạng giọng nói được đưa vào bằng chứng.

Chế độ giọng nói cho phép học sinh nhập bằng tiếng nói thay vì văn bản. Dịch vụ nhận dạng của Groq trả bản chép lời kèm độ tin cậy, còn lời gia sư được phát bằng giọng tổng hợp chạy trên máy của người vận hành. Bản chép lời độ tin cậy thấp không được coi là câu trả lời chắc chắn: độ tin cậy đi vào trọng số bằng chứng và lượt nhận dạng kém được lưu để giáo viên nghe lại. Giọng nói chỉ là phương thức nhập và xuất, không phải thành phần quyết định: sau nhận dạng, phản hồi đi qua cùng quy trình như văn bản nên mọi bất biến ở mục~\ref{sec:nfr} giữ nguyên.
